\begin{proposition}[6.15.7]
\EGAReaderTarget{ega4.c04.prop.6-15-7}{IV.2.s6.15.7-fr}
Sean $k$ un cuerpo, $X$ un $k$-preesquema, $k'$ una extensión
de $k$ y $X'=X\otimes_k k'$. Sean $x'$ un punto de $X'$ y $x$
su proyección en $X$. Para que $X$ sea geométricamente
unirramoso en el punto $x$, es necesario y suficiente que
$X'$ sea geométricamente unirramoso en el punto $x'$. Para
que $X$ sea geométricamente unirramoso, es necesario y
suficiente que $X'$ lo sea.
\end{proposition}

\begin{proof}
La segunda afirmación resulta de la primera y del hecho de
que la proyección $f:X'\to X$ es un morfismo sobreyectivo.
Para demostrar la primera afirmación, es posible limitarse, en
virtud de (\textbf I, 5.1.8), al caso en que $X$ es reducido y, en
virtud de (\textbf I, 3.6.5), al caso
\oldpage[IV-2]{180}
en que $X=\operatorname{Spec}(A)$ (con $A=\mathcal O_{X,x}$) es un
esquema \emph{local}. Observemos que $A'=\mathcal O_{X',x'}$ es, por
hipótesis, un $A$-módulo fielmente plano, que por tanto
contiene a $A$. En consecuencia, tanto la hipótesis de que
$X$ es geométricamente unirramoso en $x$ como la de que $X'$
lo es en $x'$ implican que $A$ es un anillo local \emph{íntegro}
(pues $A$, al ser reducido, es también isomorfo a un
subanillo de $A'_{\mathrm{red}}$). Sean $K$ el cuerpo de
fracciones de $A$ y $B$ la clausura integral de $A$, y
pongamos $Y=\operatorname{Spec}(B)$. Decir que $X$ es
geométricamente unirramoso en $x$ equivale a decir que el
morfismo $g:Y\to X$ es radicial en $x$ (6.15.3). Sea
$Y'=Y\otimes_k k'=Y\times_X X'$, de modo que se tiene el
diagrama conmutativo
\[
\xymatrix{
Y \ar[d]_g & Y' \ar[l]_h \ar[d]_{g'}\\
X & X' \ar[l]^f
}
\]
Observemos que $f$ es un morfismo plano; por tanto (6.15.4.1),
el morfismo entero $g'_{\mathrm{red}}$ es birracional. Por
otra parte, como $Y$ es normal, $Y'$ es geométricamente
unirramoso (6.15.6). Para que $X'$ sea geométricamente
unirramoso en $x'$, es por tanto necesario y suficiente que
$g'$ sea radicial en $x'$ (6.15.5). Ahora bien, $x$ es la
proyección de $x'$ en $X$, y decir que $g$ es radicial en $x$
equivale a decir que $g'$ es radicial en $x'$ (6.15.3.1).
Finalmente, para que $X$ sea geométricamente unirramoso en
$x$, es necesario y suficiente que $g$ sea radicial en ese
punto, de donde resulta la proposición.
\end{proof}
